\documentclass{article}
\usepackage{collection}
\usepackage{DraTex}
\usepackage{AlDraTex}
\usepackage{textcomp}
\usepackage{tikz}
\usepackage{pgffor}
\usepackage{graphicx}
\usetikzlibrary{positioning,shapes,shadows,arrows}

\begin{document}


\begin{center}
{\Huge Collection Package}
\end{center}
\indent \textbf{Objective:} is to aid in the visualization of Data Structures and Objects in the Theory of Computer Science, as well as, being flexible enough to accept any object(text,pictures,or any graphics) as input and customize to the user's desire.
\newpage
%%%CS versions
\begin{center}
\textbf{{\Large Computer Science Data Structures: Linked Lists}}
\end{center}

CS Singly Linked List\\
Format:{\textbackslash}cslinklist[Option:styling]{\{}objects{\}}\\
Example:{\textbackslash}cslinklist{\{}a,b,c,d,e{\}} \newline \\
\begin{adjustbox}{width=0.5\paperwidth,center,keepaspectratio}
\cslinklist{a,b,c,d,e}
\end{adjustbox}\\
CS Doubly Linked List\\
Format:{\textbackslash}csdlinklist[Option:styling]{\{}objects{\}}\\
Example:{\textbackslash}csdlinklist{\{}foo,boo,coo,koo,woo,hoo{\}} \newline \\
\begin{adjustbox}{width=0.5\paperwidth,center,keepaspectratio}
\csdlinklist{foo,boo,coo,koo,woo,hoo}
\end{adjustbox}\\
CS Circular Singly Linked List\\
Format:{\textbackslash}cslooplinklist[Option:styling]{\{}objects{\}}\\
Example:{\textbackslash}cslooplinklist{\{}a,b,c,d{\}} \newline \\
\begin{adjustbox}{width=0.5\paperwidth,center,keepaspectratio}
\cslooplinklist{a,b,c,d}
\end{adjustbox}\\
CS Circular Doubly Linked List\\
Format:{\textbackslash}csloopdlinklist[Option:styling]{\{}objects{\}}\\
Example:{\textbackslash}csloopdlinklist{\{}a,b,c,d{\}} \newline \\
\begin{adjustbox}{width=0.5\paperwidth,center,keepaspectratio}
\csloopdlinklist{a,b,c,d}
\end{adjustbox}\\
%%%simple form
\newpage
\begin{center}
\textbf{{\Large Simple collections: Linked Lists}}
\end{center}
Simple Singly Linked List\\
Format:{\textbackslash}linklist[Option:styling]{\{}objects{\}}\\
Example:{\textbackslash}linklist{\{}a,b,c,d{\}} \newline \\
\begin{adjustbox}{width=0.5\paperwidth,center,keepaspectratio}
\linklist{a,b,c,d}
\end{adjustbox}\\

Example:{\textbackslash}linklist{\{}{\{}{\textbackslash}tikz{\textbackslash}node[draw,circle]{\{}0{\}};{\}},{\{}{\textbackslash}tikz{\textbackslash}node[draw,circle]{\{}gh{\}};{\}},{\{}{\textbackslash}tikz{\textbackslash}node[draw,circle]{\{}b{\}};{\}},{\{}{\textbackslash}tikz{\textbackslash}node[draw,circle]{\{}c{\}};{\}}{\}}\\
\begin{adjustbox}{width=0.5\paperwidth,center,keepaspectratio}
\linklist{{\tikz\node[draw,circle]{0};},{\tikz\node[draw,circle]{gh};},{\tikz\node[draw,circle]{b};},{\tikz\node[draw,circle]{c};}} 
\end{adjustbox}\\


Simple Doubly Linked List\\
Format:{\textbackslash}dlinklist[Option:styling]{\{}objects{\}}\\
Example:{\textbackslash}dlinklist{\{}a,b,c,d{\}} \newline \\
\begin{adjustbox}{width=0.5\paperwidth,center,keepaspectratio}
\dlinklist{a,b,c,d}
\end{adjustbox}\\
Simple loop Linked List\\
Format:{\textbackslash}looplinklist[Option:styling]{\{}objects{\}}\\
Example:{\textbackslash}looplinklist{\{}a,b,c,d{\}} \newline \\
\begin{adjustbox}{width=0.5\paperwidth,center,keepaspectratio}
\looplinklist{a,b,c,d}
\end{adjustbox}\\
Simple Doubly Linked List\\
Format:{\textbackslash}loopdlinklist[Option:styling]{\{}objects{\}}\\
Example:{\textbackslash}loopdlinklist{\{}a,b,c,d{\}} \newline \\
\begin{adjustbox}{width=0.5\paperwidth,center,keepaspectratio}
\loopdlinklist{a,b,c,d}
\end{adjustbox}
\newpage
\begin{center}
\textbf{{\Large Option(Vertical): CS vs Simple: Singly and Doubly Linked Lists}}
\end{center}

\begin{center}
\begin{adjustbox}{height=0.6\paperheight,width=\paperwidth,center,keepaspectratio}
\lrmatrix{4}{{{\textbackslash}csvdlinklist{\{}foo,boo,koo{\}}},{{\textbackslash}csvlinklist{\{}a,b,c{\}}},{{\textbackslash}vlinklist{\{}a,b,c,d,e,f,g,h,i{\}}},{{\textbackslash}dvlinklist{\{}a,b,c,d,e,f,g,h,i{\}}},\csvdlinklist{foo,boo,koo},\csvlinklist{a,b,c},\vlinklist{a,b,c,d,e,f,g,h,i},\vdlinklist{a,b,c,d,e,f,g,h,i}}
\end{adjustbox}
\end{center}

%\begin{adjustbox}{width=0.8\paperwidth,center,keepaspectratio}
%\hashmapset{namew}{1,2,3}{\cslinklist[{scale=0.6,every node/.style={transform shape}}]{a,c,d},%\csdlinklist[scale=0.6, every node/.style={transform shape}]{a,c,d},hello}
%\end{adjustbox}
%%%%%%%%%%%%%%%%%%%%%%%%%%%%%%%%%%%%%%%%%%%
%%%%%%%%%%%%%%%%%%%%%%%%%%%%%%%%%%%%%%


%%%%%%%%%%%%%%%%%%%%%%%%%%%%%%%%%%%%%%%%%%%%%%%%
\newpage
\begin{center}
\textbf{{\Large Data Structure: HashMap Object}}
\end{center}
Format:{\textbackslash}hashmapset[Option:styling]{\{}Name of Hashmap{\} }{\{}list of Keys{\}}{\{}list of Values{\}}\\
Example:{\textbackslash}hashmapset{\{}HashMap{\} }{\{}key1,key2,key3{\}}{\{}value1,value2,value3{\}}\newline \\
\begin{adjustbox}{width=0.5\paperwidth,center,keepaspectratio}
\hashmapset{HashMap}{key1,key2,key3}{value1,value2,value3}
\end{adjustbox}

\newpage
\begin{center}
\textbf{{\Large Data Structure: Stacks}}
\end{center}

\section*{Vertical Stack}
{\textbackslash}begin{\{}stack{\}}{\{}vertical{\}}
  {\textbackslash}cell[blue]{\{}cell 1{\}}{\{}1{\}}
{\textbackslash}end{\{}stack{\}} \\ \newline
\begin{stack}{vertical}
  \cell[blue]{cell 1}{1}
\end{stack}\\ \newline

\section*{Horizontal Stack}
{\textbackslash}begin{\{}stack{\}}{\{}horizontal{\}}
  {\textbackslash}cell[blue]{\{}cell 1{\}}{\{}1{\}}
{\textbackslash}end{\{}stack{\}} \\ \newline
\begin{stack}{horizontal}
  \cell[blue]{cell 1}{1}
  \cell[red]{cell 2}{2}
\end{stack}

\newpage
\begin{center}
\textbf{{\Large ARRAYS:}}
\end{center}
(Left-Right)\\ 
Format:{\textbackslash}lrmatrix{\{}number of columns to split list of objects{\} }{\{}objects{\}}\\
Example1:{\textbackslash}lrmatrix{\{}4{\} }{\{}1,2,3,4,5,6,7,a,b,c{\}}\\
Example2:{\textbackslash}lrmatrix{\{}2{\} }{\{}{\textbackslash}includegraphics{\{}ex1.jpg{\}},{\textbackslash}includegraphics{\{}ex2.jpg{\}},\\{\textbackslash}includegraphics{\{}ex3.jpg{\}},{\textbackslash}includegraphics{\{}ex4.jpg{\}}{\}}\newline \\

\lrmatrix{4}{1,2,3,4,5,6,7,a,b,c} \lrmatrix{2}{\includegraphics[width=30mm]{computer_science_zoomed_by_blackblood357-d37ofou.jpg},\includegraphics[width=30mm]{computer-science1.jpg},\includegraphics[width=30mm]{computer-science1.jpg},\includegraphics[width=30mm]{computer_science_zoomed_by_blackblood357-d37ofou.jpg}} \newline \\

(Vertical)\\
Format:{\textbackslash}vmatrix{\{}number of rows to split list of objects{\}}{\{}objects{\}}\\
Example:{\textbackslash}vmatrix{\{}2{\}}{\{}1,2,3,4,5,6,7,a,b,c{\}}\\
\vmatrix{2}{1,2,3,4,5,6,7,a,b,c}
\newpage
\textbf{Data Structure:(Simple Hashtable)}\\
{\textbackslash}lrmatrix{\{}1{\} }{\{}\\{\textbackslash}linklist{\{}{\{}{\textbackslash}tikz{\textbackslash}node[draw,circle]{\{}0{\}};{\}},a,b,c{\}},\\{\textbackslash}linklist{\{}{\{}{\textbackslash}tikz{\textbackslash}node[draw,circle]{\{}1{\}};{\}},d,e,f{\}},\\{\textbackslash}linklist{\{}{\{}{\textbackslash}tikz{\textbackslash}node[draw,circle]{\{}0{\}};{\}},g,h,i{\}}{\}}\newline \\
\lrmatrix{1}{\linklist{{\tikz\node[draw,circle]{0};},a,b,c},\linklist{{\tikz\node[draw,circle]{1};},d,e,f},\linklist{{\tikz\node[draw,circle]{2};},g,h,i}}


\section*{Diagrams}
Format:\\
{\textbackslash}begin{\{}diagram{\}}{\{}Styling{\}}\\ 
-----Content------\\
{\textbackslash}end{\{}diagram{\}}\\
\subsection*{Example: Class Diagram}
Format:\\
{\textbackslash}class[Option: styling]{\{}class var{\}}[Option:Title][Option:members][Option:operations]\\
{\textbackslash}rel{\{}relationship{\}}{\{}derive source{\}} {\{}deriver{\}}\\
EX:
{\textbackslash}begin{\{}diagram{\}}{\{}node distance=5cm,layered layout{\}}\\
{\textbackslash}class[rounded corners, fill=red!20]{\{}a{\}}[Interface{\textbackslash}{\textbackslash}Classes][members][operations]\\
{\textbackslash}class[rounded corners, fill=blue!20]{\{}b{\}}[class1][members][operations]\\
{\textbackslash}class{\{}c{\}}[class2][members][operations]\\
{\textbackslash}class{\{}d{\}}[class3][members]\\
{\textbackslash}class[rounded corners]{\{}e{\}}[other Object][members][operations]\\
{\textbackslash}rel{\{}generalization{\}}{\{}a{\}} {\{}b,c,d{\}}\\
{\textbackslash}rel{\{}dependency{\}}{\{}e{\}}{\{}d{\}}\\  
{\textbackslash}end{\{}diagram{\}}\\

\begin{diagram}{node distance=5cm,layered layout}
\class[rounded corners, fill=red!20]{a}[Interface\\Classes][members][operations]
\class[rounded corners, fill=blue!20]{b}[class1][members][operations]
\class{c}[class2][members][operations]
\class{d}[class3][members]
\class[rounded corners]{e}[other Object][members][operations]
\rel{generalization}{a} {b,c,d}
\rel{dependency}{e}{d}
    
\end{diagram}

%%%%%%%1 input
\newpage
\begin{center}
\textbf{{\Large Objects:}}
\end{center}
Example1:{{\textbackslash}numofnodes{1}{\\}}\\
Example2:{{\textbackslash}numofnodes{2}{\\}}\\
Example3:{{\textbackslash}numofnodes{3}{\\}}\\
\newpage
\begin{center}
\textbf{{\Large Objects:}}
\end{center}
Example1:{{\textbackslash}numofnodes{1}{\\}}\\
\tikzstyle{abstract}=[rectangle, draw=black, fill=white!40,
        text centered, anchor=north, text=black, text width=4cm]
\usetikzlibrary{arrows,positioning,shapes}
\newcommand{\numofnodes}{1}

%User Input
%%Change the "Caption Here" to the caption
\newcommand{\Caption}{Caption Here}

%Head
\newcommand{\ObjectName}{LinkedList\textlangle student\textrangle}
\newcommand{\LLHead}{}
\newcommand{\LLTail}{null}
\newcommand{\LLSize}{\numofnodes}

%Node
\newcommand{\NName}{Node\textlangle student\textrangle}
\newcommand{\NValue}{}
\newcommand{\NNext}{}
\newcommand{\NPrev}{null}

%Value
\newcommand{\VTop}{Student}
\newcommand{\VName}{Jim}
\newcommand{\VNum}{123}
\newcommand{\VAddr}{}

%Address
\newcommand{\AName}{Address}
\newcommand{\AStreet}{33 Blank St}
\newcommand{\ACity}{Glassboro}
\newcommand{\AState}{NJ}
\newcommand{\AZip}{08080}

%Linked List Tail
\newcommand{\TName}{Tail}
\newcommand{\THead}{null}
\newcommand{\TTail}{}
\newcommand{\TSize}{null}

\begin {figure}

%Head
%%Draws the head node
\Draw


\MinNodeSize (20,10)
\ArrowHeads(1)
\ArrowSpec(F,10,4,5)
%% \DecVar \MaxX

%% Display the primitive values in an object
%% Called once for each primitive field
\Define \Var(3)
{	
%% parms:
%%    #1: 0 if this var has no primitive value
%%    #2: name of the var
%%    #3: value in a box
	
   \IF \GtInt (#1,0)  \THEN
	\MarkLoc(left)
 	\Move(0,-27)
	\EntryExit(-1,0,-1,0)
	\Node (tag) (--#2--)
	%% Note the longest field value, to draw the box
	{ 
	\MoveToExit(1,0)
	  \ORectNode (tag) (--#3--)
	\MoveToExit(1,0)  \MarkLoc(right)
	\Move(50,0)  \MarkLoc(rightPlus) 
	\MoveToLL(right,rightPlus)(topL,maxL)
	\MarkGLoc(bottomL)  
	\LSeg\Q(maxL,maxRight)
	\LSeg\R(bottomL,right)
	\IF \GtDec (\R,\Q) \THEN
	    \MoveToLoc(right)
	    \MarkGLoc(maxRight)
	    \FI
	}
	\EntryExit(0,0,0,0)
\MarkGLoc(bottom)
   \FI
}


%% Display the object references in an object
%% Called once for each primitive field
%% Object is a recursive data structure
\Define \Obj(8)
{	
%% parms:
%%    #1: 0 do not display
%%        1 display ordinary ref
%%	  2 display ref in a set
%%        3 display a key in a map
%%        4 display a ref to a set
%%        5 display a ref to a map
%%  	  6 display object only, calling OD directly, no arrow
%%    #2: name of the field
%%    #3: class of the field  (for object types)
%%    #4: table of primitive fields in the referred object
%% 	   or size for a map
%%    #5: table of object fields in the referred object
%%         or table of primitive keys for maps
%%    #6: table of object keys in the referred object
%%    #7: table of primitive values in the referred object
%%    #8: table of object values in the referred object

  % Base case for recursion:  #1 = 0

\IF \GtInt (#1, 0) \THEN		% Recursive case
					% for ordinary field or value in map

   \Move(0,-25)
   \MarkLoc(tmp)
   \EntryExit(-1,0,1.1,0)
   \Node (tag) (--#2--)
   \ORectNode(ref) (-- --)
   \EntryExit(0,0,0,0)
   \IF \EqInt(#1,1) \THEN		% ordinary ref
      \Move(100,\Val\I)
      \OD(#3, #4, #5, 1)  
   \FI
   \IF \EqInt(#1,2) \THEN		% ref in a set
 	\ODset(#3,#4,#5)
   \FI
   \IF \EqInt(#1,3) \THEN		% key in a map
	    \Move(-160,\Val\I)
	    \Move(0,-10)    
	    \OD(#3, #4, #5,3)  
   \FI
   \IF \EqInt(#1,4) \THEN		% ref to a set
	       \OD(#3,#4,#5,4)
   \FI
   \IF \EqInt(#1,5) \THEN		% ref to a map
	          \ODmap(#3,#4,#5,#6,#7,#8,#1)
   \FI
   \MoveToLoc(tmp)
   \I-100;
   \MarkGLoc(bottom)
\FI
   }
 
%% Draw the object in a rectangular box
%% Show primitive values
%% Show references as arrows to objects
\Define \OD(4)
{  
%% parms:
%%    #1:  class of the object 
%%    #2:  table of primitive fields 
%%    #3:  table of object fields 
%%    #4:  See parm #1 for Obj above

   \SaveAll
%%%%  J is depth of recursion
\J+1;
\IF \GtInt (\J,2) \THEN
      \Move(-400,-200)
      \J=0;
    
\FI
\Move (-20,60)   %%% ???  
\Node(obj)(--\underline{#1}--)
%% topLeft and bottomRight are used to draw the rectangle
\MoveToExit(1.02,0)    \MarkGLoc(maxRight)   %% maxRight stores largest primitive
\MoveToExit(-1.2,1.2)  \MarkLoc(topLeft)
\Move (50,-70)
 
\MinNodeSize(30,15)
\PenSize(0.25pt)	%% Lightweight arrow, might cross over other objects

%% Draw the arrow for a reference
\IF \EqInt (#4,3) \THEN 	%% key in a map
     { \CurvedEdgeAt (ref,-0.8,0, obj,1,0.5)(180,0.6,90, 0.6) }
	\ELSE
	\IF \EqInt (#4,6) \THEN		%% calling OD directly 
		{ } 			%% no arrow
        \ELSE
     	{ \CurvedEdgeAt (ref,0.8,0, obj,-1,0.8)(0,0.6,180, 0.6) }
   	\FI
   \FI

\Move(-45,70)
\PenSize(0.75pt)

{\MarkGLoc (maxL) \Move (0,100) \MarkGLoc(topL) }

\Indirect <#2>(0,9) { \Var }		%% Draw primitive fields
\MarkLoc(bottomTemp)
\MoveToLoc(maxRight) \MarkLoc(maxRightLocal)
\MoveToLoc(bottomTemp)

\Indirect<#3>(0,8) { \Obj   }		%% Draw object fields

%% Set maxRightLocal in case there were no primitive fields
\Move(100,0) \MarkLoc(maxRightTemp)
\MoveToRightmost(maxRightLocal,maxRightTemp)
\MarkLoc(maxRightLocal)
{
%% Find bottom right corner of rectangle
\MoveToLoc(bottom)
\Move(100,0) \MarkLoc(bottomR)
\MoveToLoc(maxRightLocal) \Move(0,-100) \MarkLoc(rightB)
\MoveToLL(bottom,bottomR)(maxRightLocal,rightB)
\Move(10,-20) 
\MarkLoc(bottomRight)
\CSeg \DrawRect(bottomRight,topLeft)
}
   
   \RecallAll
\J-1;
}

%% Draw a ref to a Set
\Define \ObjSet(6)
{	
%% parms:
%%    #1: 0 iff this field has no Object ref
%%    #2: name of the field
%%    #3: class of the field  (for object types)
%%    #4: table of primitive fields in the referred object - wrappers or Strings
%%    #5: table of object fields in the referred object
%%    #6: size of the set 

  % Base case:  #1 = 0

\IF \GtInt (#1, 0) \THEN		% Recursive case
	\MarkLoc(tmp)
   	\Node (tag) (--#2--)
	\Move (50,0)
	\ORectNode(ref) (-- --)
	\Move(50,\Val\I)
	\Move (35,20)
	\ODset(#3, #4, #5,#6)  
	\MoveToLoc(tmp)
	\I-135;
   \FI
}

%% Draw a set with the refs at 'random' places in the box
\Define \ODset(4)    
{  
%% parms:
%%    #1:  class of the object
%%    #2:  table of primitive fields 
%%    #3:  table of object fields 
%%    #4:  size of the set        
   \SaveAll
  
\MinNodeSize (90,110)
\RectNode(obj)(-- #1~~~~~~~~~~~~~~~~~~~~~~~--)
\MinNodeSize (20,10)

{ \PenSize(0.25pt)	%% Lightweight arrow, might cross over other objects
  %% Draw arrow to this Set object
  \CurvedEdgeAt (ref,0.8,0, obj,-1,0.8)(10,0.6,180, 0.6) }
  \PenSize(0.75pt)
   \IntVar \X      
   \IntVar \Y

   \Move(-50,70)
   { \Line(110,0) }
   \Move (15,-15)
   \Node(tag)(--size--)
   \MoveToExit(1,0)
   \Move(25,0)
   \ORectNode(tag)(--#4--)
   \Move(-90,0)
   \MarkLoc(start)

%% Place the fields at seemingly random locations
   \Indirect <#2>(0,9) { 		%% Draw the primitive fields
 	\X+34;
 	\K=\X;
	\K/90;
	\K*90;
 	\X-\K;		%% x + 34 mod 90
	\Y+73;
	\K=\Y;
	\K/110;
	\K*110;
	\Y-\K;		%% y + 73 mod 110
	\Y+20;		%% Allow for size at top
 		\MoveToLoc(start)
		\Move (\Val\X,-\Val\Y)
 		\Var
}

%% Place the fields at seemingly random locations
   \Indirect<#3>(0,9) { 		%% Draw the object fields
 	\X+34;
 	\K=\X;
 	\K/90;
 	\K*90;
 	\X-\K;		%% x + 72 mod 90
	\Y+73;
	\K=\Y;
	\K/110;
	\K*110;
	\Y-\K;		%% y + 73 mod 110
	\Y+20;		%% Allow for size at top
 	\MoveToLoc(start)
	\Move (\Val\X,-\Val\Y)
		 \Obj   
	}
   \RecallAll
}

%% Draw a map object in a rectangle
%% Value refs point to the right
%% Key refs point to the left
%% Note largest value field to draw the rectangle
\Define \ODmap(7)
{  
%% parms:
%%    #1:  class of the object 
%%    #2:  size
%%    #3:  table of primitive keys 
%%    #4:  table of object keys 
%%    #5:  table of primitive values 
%%    #6:  table of object values 
%%    #7:  see parm #1 for Obj above
   \SaveAll
  
\Move(-20,-60) 
\Node(obj)(--\underline{#1}--)
\MoveToExit(1.02,0) 	\MarkLoc (maxRightLocal)
			\MarkLoc (maxRight)
\MoveToExit(-1.2,1.2)\MarkLoc (topLeft)

\MinNodeSize(30,15)
\PenSize(0.25pt)		%% Use Lightweight pen, may cross other objects

\IF \EqInt(#7,3) \THEN
     \CurvedEdgeAt (ref,-0.8,0, obj,1,0.5)(180,0.6,90, 0.6) 
     \ELSE
	\IF \EqInt(#7,6) \THEN		%% calling ODmap directly
	    {  }
	\ELSE
	  \CurvedEdgeAt (ref,0.8,0, obj,-1,0.8)(10,0.6,180, 0.6) 
        \FI
\FI
	
\PenSize(0.75pt)
\Move(20,-35)
   
   %% Show size and headings for keys and values
   { \Node(tag)(--size--) \Move(40,0) \ORectNode(tag)(--#2--) }
   \Move (10,-20)
   { \Node(tag)(--\underline{keys}--) 
     \Move(75,0) \Node(tag)(--\underline {values}--) 
   }
   \Move (-50,0)
   \MarkLoc(keys)

   {\MarkGLoc (maxL) \Move (0,100) \MarkGLoc(topL) }   %% ??

  %% keys 
   \Indirect <#3>(0,9) { \Var  }	%% Draw the primitive keys

   \Indirect<#4>(0,8) { \Obj   }	%% Draw the object keys
   \MarkGLoc(bottom)
\MoveToLoc(keys)
\Move(83, -1)
{
  %% values 
   \Indirect <#5>(0,9) { \Var  }	%% Draw the primitive values
   \MarkLoc(bottomTemp)
   \MoveToLoc(bottomTemp)

   \Indirect<#6>(0,8) { \Obj   }	%% Draw the object values

%% Keys may not be primitive, affects size of box
\Move(100,0)	\MarkLoc(maxRightTemp)
\MoveToRightmost(maxRightLocal,maxRightTemp)
\MarkLoc(maxRightLocal)
}
{
%% Find bottom right corner of rectangle
\MoveToLoc(bottom)
\Move(100,0)	\MarkLoc(bottomR)
\MoveToLoc(maxRightLocal)	\Move(0,-100)	\MarkLoc(rightB)
\MoveToLL(bottom,bottomR)(maxRightLocal,rightB)
\Move(10,-20)
\MarkLoc(bottomRight)
\CSeg \DrawRect(bottomRight,topLeft)
}
   \RecallAll
}

%% Move the pen to the loc which is furthest right
\Define \MoveToRightmost(2)
{
%% #1:  A loc
%% #2:  A loc
\MoveTo(-200,100)	\MarkLoc(upLeft)
\MoveTo(-200,-100)	\MarkLoc(downLeft)
\MoveToLoc(#1)		\Move(100,0)	\MarkLoc(right)
\MoveToLL(upLeft,downLeft)(#1,right)	\MarkLoc(leftA)
\MoveToLoc(#2)		\Move(100,0)	\MarkLoc(right)
\MoveToLL(upLeft,downLeft)(#2,right)	\MarkLoc(leftB)

	\LSeg\Q(leftA,#1)
	\LSeg\R(leftB,#2)
	\MoveToLoc(#1)
	\IF \GtDec (\R,\Q) \THEN
	    \MoveToLoc(#2)
	    \FI
}

\hspace{-4cm}

\Indirect \Table <ObjectPrims>
{  
	1, head, \LLHead  &
    1, tail, \LLTail &
    1, size, \LLSize
}

\Indirect \Table <ObjectDiagram>
 {  0, home,Address,addrAPrims,addrAObjs, , , }
\Scale (1,1)
\OD (\ObjectName, ObjectPrims, ObjectDiagram, 6)

%Arrow from head
%%Draws the arrow from head to first node
\Move(5,43)
\FcNode(Head)
\Move(45,0)
\FcNode(HeadRight)
\Move(0,-80)
\FcNode(HeadDown)
\Move(-100,0)
\FcNode(HeadLeft)
\Move(0,-14)
\FcNode(HeadBottom)
\ArrowHeads(1)
\Edge(HeadLeft,HeadBottom)
\ArrowHeads(0)
\Edge(Head,HeadRight)
\VHEdge(HeadRight,HeadDown)
\Edge(HeadDown,HeadLeft)
\EndDraw

%Node
\begin{tikzpicture}[abstract, rectangle split, rectangle split parts=2, text justified]
\foreach \x in {0}
    \foreach \y in {1,...,\numofnodes} 
       {\node (NODES) [abstract, rectangle split, rectangle split parts=2, text justified]  (\x\y) at (5, 2.9*\y) 
       		{\textbf{Node}
            \nodepart{second}next  \framebox[1.1\width]{arrow}
            \newline prev  \framebox[1.1\width]{arrow}
            \newline value  \framebox[1.1\width]{arrow}
            
            };}
                      
\foreach \x in {0}
    \foreach \y in {1,...,\numofnodes}            
{\node (NODETHING) [abstract, rectangle split, rectangle split parts=2, text justified] (\x\y) at (10, 2.6*\y) 
       		{\textbf{\VTop}
            \nodepart{second}Name  \framebox[1.1\width]{\VName}
            \newline Num  \framebox[1.1\width]{\VNum}
            \newline Address  \framebox[1.1\width]{arrow}};}

\foreach \x in {0}
    \foreach \y in {1,...,\numofnodes} 
{\node (NODETHING2) [abstract, rectangle split, rectangle split parts=2, text justified] (\x\y) at (15, 2.6*\y) 
       		{\textbf{\AName}
            \nodepart{second}Street  \framebox[1.1\width]{\AStreet}
            \newline City  \framebox[1.1\width]{\ACity}
            \newline State  \framebox[1.1\width]{\AState}
            \newline Zip  \framebox[1.1\width]{\AZip}};}           
\end{tikzpicture}

%%Tail
%%Draws the tail node
\hspace{.05cm}
\Draw


\MinNodeSize (20,10)
\ArrowHeads(1)
\ArrowSpec(F,10,4,5)
%% \DecVar \MaxX

%% Display the primitive values in an object
%% Called once for each primitive field
\Define \Var(3)
{	
%% parms:
%%    #1: 0 if this var has no primitive value
%%    #2: name of the var
%%    #3: value in a box
	
   \IF \GtInt (#1,0)  \THEN
	\MarkLoc(left)
 	\Move(0,-27)
	\EntryExit(-1,0,-1,0)
	\Node (tag) (--#2--)
	%% Note the longest field value, to draw the box
	{ 
	\MoveToExit(1,0)
	  \ORectNode (tag) (--#3--)
	\MoveToExit(1,0)  \MarkLoc(right)
	\Move(50,0)  \MarkLoc(rightPlus) 
	\MoveToLL(right,rightPlus)(topL,maxL)
	\MarkGLoc(bottomL)  
	\LSeg\Q(maxL,maxRight)
	\LSeg\R(bottomL,right)
	\IF \GtDec (\R,\Q) \THEN
	    \MoveToLoc(right)
	    \MarkGLoc(maxRight)
	    \FI
	}
	\EntryExit(0,0,0,0)
\MarkGLoc(bottom)
   \FI
}


%% Display the object references in an object
%% Called once for each primitive field
%% Object is a recursive data structure
\Define \Obj(8)
{	
%% parms:
%%    #1: 0 do not display
%%        1 display ordinary ref
%%	  2 display ref in a set
%%        3 display a key in a map
%%        4 display a ref to a set
%%        5 display a ref to a map
%%  	  6 display object only, calling OD directly, no arrow
%%    #2: name of the field
%%    #3: class of the field  (for object types)
%%    #4: table of primitive fields in the referred object
%% 	   or size for a map
%%    #5: table of object fields in the referred object
%%         or table of primitive keys for maps
%%    #6: table of object keys in the referred object
%%    #7: table of primitive values in the referred object
%%    #8: table of object values in the referred object

  % Base case for recursion:  #1 = 0

\IF \GtInt (#1, 0) \THEN		% Recursive case
					% for ordinary field or value in map

   \Move(0,-25)
   \MarkLoc(tmp)
   \EntryExit(-1,0,1.1,0)
   \Node (tag) (--#2--)
   \ORectNode(ref) (-- --)
   \EntryExit(0,0,0,0)
   \IF \EqInt(#1,1) \THEN		% ordinary ref
      \Move(100,\Val\I)
      \OD(#3, #4, #5, 1)  
   \FI
   \IF \EqInt(#1,2) \THEN		% ref in a set
 	\ODset(#3,#4,#5)
   \FI
   \IF \EqInt(#1,3) \THEN		% key in a map
	    \Move(-160,\Val\I)
	    \Move(0,-10)    
	    \OD(#3, #4, #5,3)  
   \FI
   \IF \EqInt(#1,4) \THEN		% ref to a set
	       \OD(#3,#4,#5,4)
   \FI
   \IF \EqInt(#1,5) \THEN		% ref to a map
	          \ODmap(#3,#4,#5,#6,#7,#8,#1)
   \FI
   \MoveToLoc(tmp)
   \I-100;
   \MarkGLoc(bottom)
\FI
   }
 
%% Draw the object in a rectangular box
%% Show primitive values
%% Show references as arrows to objects
\Define \OD(4)
{  
%% parms:
%%    #1:  class of the object 
%%    #2:  table of primitive fields 
%%    #3:  table of object fields 
%%    #4:  See parm #1 for Obj above

   \SaveAll
%%%%  J is depth of recursion
\J+1;
\IF \GtInt (\J,2) \THEN
      \Move(-400,-200)
      \J=0;
    
\FI
\Move (-20,60)   %%% ???  
\Node(obj)(--\underline{#1}--)
%% topLeft and bottomRight are used to draw the rectangle
\MoveToExit(1.02,0)    \MarkGLoc(maxRight)   %% maxRight stores largest primitive
\MoveToExit(-1.2,1.2)  \MarkLoc(topLeft)
\Move (50,-70)
 
\MinNodeSize(30,15)
\PenSize(0.25pt)	%% Lightweight arrow, might cross over other objects

%% Draw the arrow for a reference
\IF \EqInt (#4,3) \THEN 	%% key in a map
     { \CurvedEdgeAt (ref,-0.8,0, obj,1,0.5)(180,0.6,90, 0.6) }
	\ELSE
	\IF \EqInt (#4,6) \THEN		%% calling OD directly 
		{ } 			%% no arrow
        \ELSE
     	{ \CurvedEdgeAt (ref,0.8,0, obj,-1,0.8)(0,0.6,180, 0.6) }
   	\FI
   \FI

\Move(-45,70)
\PenSize(0.75pt)

{\MarkGLoc (maxL) \Move (0,100) \MarkGLoc(topL) }

\Indirect <#2>(0,9) { \Var }		%% Draw primitive fields
\MarkLoc(bottomTemp)
\MoveToLoc(maxRight) \MarkLoc(maxRightLocal)
\MoveToLoc(bottomTemp)

\Indirect<#3>(0,8) { \Obj   }		%% Draw object fields

%% Set maxRightLocal in case there were no primitive fields
\Move(100,0) \MarkLoc(maxRightTemp)
\MoveToRightmost(maxRightLocal,maxRightTemp)
\MarkLoc(maxRightLocal)
{
%% Find bottom right corner of rectangle
\MoveToLoc(bottom)
\Move(100,0) \MarkLoc(bottomR)
\MoveToLoc(maxRightLocal) \Move(0,-100) \MarkLoc(rightB)
\MoveToLL(bottom,bottomR)(maxRightLocal,rightB)
\Move(10,-20) 
\MarkLoc(bottomRight)
\CSeg \DrawRect(bottomRight,topLeft)
}
   
   \RecallAll
\J-1;
}

%% Draw a ref to a Set
\Define \ObjSet(6)
{	
%% parms:
%%    #1: 0 iff this field has no Object ref
%%    #2: name of the field
%%    #3: class of the field  (for object types)
%%    #4: table of primitive fields in the referred object - wrappers or Strings
%%    #5: table of object fields in the referred object
%%    #6: size of the set 

  % Base case:  #1 = 0

\IF \GtInt (#1, 0) \THEN		% Recursive case
	\MarkLoc(tmp)
   	\Node (tag) (--#2--)
	\Move (50,0)
	\ORectNode(ref) (-- --)
	\Move(50,\Val\I)
	\Move (35,20)
	\ODset(#3, #4, #5,#6)  
	\MoveToLoc(tmp)
	\I-135;
   \FI
}

%% Draw a set with the refs at 'random' places in the box
\Define \ODset(4)    
{  
%% parms:
%%    #1:  class of the object
%%    #2:  table of primitive fields 
%%    #3:  table of object fields 
%%    #4:  size of the set        
   \SaveAll
  
\MinNodeSize (90,110)
\RectNode(obj)(-- #1~~~~~~~~~~~~~~~~~~~~~~~--)
\MinNodeSize (20,10)

{ \PenSize(0.25pt)	%% Lightweight arrow, might cross over other objects
  %% Draw arrow to this Set object
  \CurvedEdgeAt (ref,0.8,0, obj,-1,0.8)(10,0.6,180, 0.6) }
  \PenSize(0.75pt)
   \IntVar \X      
   \IntVar \Y

   \Move(-50,70)
   { \Line(110,0) }
   \Move (15,-15)
   \Node(tag)(--size--)
   \MoveToExit(1,0)
   \Move(25,0)
   \ORectNode(tag)(--#4--)
   \Move(-90,0)
   \MarkLoc(start)

%% Place the fields at seemingly random locations
   \Indirect <#2>(0,9) { 		%% Draw the primitive fields
 	\X+34;
 	\K=\X;
	\K/90;
	\K*90;
 	\X-\K;		%% x + 34 mod 90
	\Y+73;
	\K=\Y;
	\K/110;
	\K*110;
	\Y-\K;		%% y + 73 mod 110
	\Y+20;		%% Allow for size at top
 		\MoveToLoc(start)
		\Move (\Val\X,-\Val\Y)
 		\Var
}

%% Place the fields at seemingly random locations
   \Indirect<#3>(0,9) { 		%% Draw the object fields
 	\X+34;
 	\K=\X;
 	\K/90;
 	\K*90;
 	\X-\K;		%% x + 72 mod 90
	\Y+73;
	\K=\Y;
	\K/110;
	\K*110;
	\Y-\K;		%% y + 73 mod 110
	\Y+20;		%% Allow for size at top
 	\MoveToLoc(start)
	\Move (\Val\X,-\Val\Y)
		 \Obj   
	}
   \RecallAll
}

%% Draw a map object in a rectangle
%% Value refs point to the right
%% Key refs point to the left
%% Note largest value field to draw the rectangle
\Define \ODmap(7)
{  
%% parms:
%%    #1:  class of the object 
%%    #2:  size
%%    #3:  table of primitive keys 
%%    #4:  table of object keys 
%%    #5:  table of primitive values 
%%    #6:  table of object values 
%%    #7:  see parm #1 for Obj above
   \SaveAll
  
\Move(-20,-60) 
\Node(obj)(--\underline{#1}--)
\MoveToExit(1.02,0) 	\MarkLoc (maxRightLocal)
			\MarkLoc (maxRight)
\MoveToExit(-1.2,1.2)\MarkLoc (topLeft)

\MinNodeSize(30,15)
\PenSize(0.25pt)		%% Use Lightweight pen, may cross other objects

\IF \EqInt(#7,3) \THEN
     \CurvedEdgeAt (ref,-0.8,0, obj,1,0.5)(180,0.6,90, 0.6) 
     \ELSE
	\IF \EqInt(#7,6) \THEN		%% calling ODmap directly
	    {  }
	\ELSE
	  \CurvedEdgeAt (ref,0.8,0, obj,-1,0.8)(10,0.6,180, 0.6) 
        \FI
\FI
	
\PenSize(0.75pt)
\Move(20,-35)
   
   %% Show size and headings for keys and values
   { \Node(tag)(--size--) \Move(40,0) \ORectNode(tag)(--#2--) }
   \Move (10,-20)
   { \Node(tag)(--\underline{keys}--) 
     \Move(75,0) \Node(tag)(--\underline {values}--) 
   }
   \Move (-50,0)
   \MarkLoc(keys)

   {\MarkGLoc (maxL) \Move (0,100) \MarkGLoc(topL) }   %% ??

  %% keys 
   \Indirect <#3>(0,9) { \Var  }	%% Draw the primitive keys

   \Indirect<#4>(0,8) { \Obj   }	%% Draw the object keys
   \MarkGLoc(bottom)
\MoveToLoc(keys)
\Move(83, -1)
{
  %% values 
   \Indirect <#5>(0,9) { \Var  }	%% Draw the primitive values
   \MarkLoc(bottomTemp)
   \MoveToLoc(bottomTemp)

   \Indirect<#6>(0,8) { \Obj   }	%% Draw the object values

%% Keys may not be primitive, affects size of box
\Move(100,0)	\MarkLoc(maxRightTemp)
\MoveToRightmost(maxRightLocal,maxRightTemp)
\MarkLoc(maxRightLocal)
}
{
%% Find bottom right corner of rectangle
\MoveToLoc(bottom)
\Move(100,0)	\MarkLoc(bottomR)
\MoveToLoc(maxRightLocal)	\Move(0,-100)	\MarkLoc(rightB)
\MoveToLL(bottom,bottomR)(maxRightLocal,rightB)
\Move(10,-20)
\MarkLoc(bottomRight)
\CSeg \DrawRect(bottomRight,topLeft)
}
   \RecallAll
}

%% Move the pen to the loc which is furthest right
\Define \MoveToRightmost(2)
{
%% #1:  A loc
%% #2:  A loc
\MoveTo(-200,100)	\MarkLoc(upLeft)
\MoveTo(-200,-100)	\MarkLoc(downLeft)
\MoveToLoc(#1)		\Move(100,0)	\MarkLoc(right)
\MoveToLL(upLeft,downLeft)(#1,right)	\MarkLoc(leftA)
\MoveToLoc(#2)		\Move(100,0)	\MarkLoc(right)
\MoveToLL(upLeft,downLeft)(#2,right)	\MarkLoc(leftB)

	\LSeg\Q(leftA,#1)
	\LSeg\R(leftB,#2)
	\MoveToLoc(#1)
	\IF \GtDec (\R,\Q) \THEN
	    \MoveToLoc(#2)
	    \FI
}


\Indirect \Table <ObjectPrims>
{  
	1, next, \THead &
    1, prev, \TTail &
    1, value, \TSize
}
\Indirect \Table <ObjectDiagram>
 {  0, home,Address,addrAPrims,addrAObjs, , , }
\Scale (1,1)
\OD (\TName, ObjectPrims, ObjectDiagram, 6)

%Arrow from Tail
%%Draws the arrow from tail node to tail
\Move(48,15)
\FcNode(Tail)
\Move(45,0)
\FcNode(Temp2)
\Move(0,60)
\FcNode(Temp3)
\Move(-100,0)
\FcNode(Temp4)
\Move(0,14)
\FcNode(Temp5)
\ArrowHeads(1)
\Edge(Temp4,Temp5)
\ArrowHeads(0)
\Edge(Tail,Temp2)
\Edge(Temp2,Temp3)
\Edge(Temp3,Temp4)
\EndDraw

\caption {\Caption}

\label {fig:ObjectDiagram}

\end {figure}


%%%%%%%2 input
\newpage
\tikzstyle{abstract}=[rectangle, draw=black, fill=white!40,
        text centered, anchor=north, text=black, text width=4cm]
\usetikzlibrary{arrows,positioning,shapes}
\newcommand{\numofnodess}{2}

\begin {figure}

%Head
%%Draws the head node
\Draw


\MinNodeSize (20,10)
\ArrowHeads(1)
\ArrowSpec(F,10,4,5)
%% \DecVar \MaxX

%% Display the primitive values in an object
%% Called once for each primitive field
\Define \Var(3)
{	
%% parms:
%%    #1: 0 if this var has no primitive value
%%    #2: name of the var
%%    #3: value in a box
	
   \IF \GtInt (#1,0)  \THEN
	\MarkLoc(left)
 	\Move(0,-27)
	\EntryExit(-1,0,-1,0)
	\Node (tag) (--#2--)
	%% Note the longest field value, to draw the box
	{ 
	\MoveToExit(1,0)
	  \ORectNode (tag) (--#3--)
	\MoveToExit(1,0)  \MarkLoc(right)
	\Move(50,0)  \MarkLoc(rightPlus) 
	\MoveToLL(right,rightPlus)(topL,maxL)
	\MarkGLoc(bottomL)  
	\LSeg\Q(maxL,maxRight)
	\LSeg\R(bottomL,right)
	\IF \GtDec (\R,\Q) \THEN
	    \MoveToLoc(right)
	    \MarkGLoc(maxRight)
	    \FI
	}
	\EntryExit(0,0,0,0)
\MarkGLoc(bottom)
   \FI
}


%% Display the object references in an object
%% Called once for each primitive field
%% Object is a recursive data structure
\Define \Obj(8)
{	
%% parms:
%%    #1: 0 do not display
%%        1 display ordinary ref
%%	  2 display ref in a set
%%        3 display a key in a map
%%        4 display a ref to a set
%%        5 display a ref to a map
%%  	  6 display object only, calling OD directly, no arrow
%%    #2: name of the field
%%    #3: class of the field  (for object types)
%%    #4: table of primitive fields in the referred object
%% 	   or size for a map
%%    #5: table of object fields in the referred object
%%         or table of primitive keys for maps
%%    #6: table of object keys in the referred object
%%    #7: table of primitive values in the referred object
%%    #8: table of object values in the referred object

  % Base case for recursion:  #1 = 0

\IF \GtInt (#1, 0) \THEN		% Recursive case
					% for ordinary field or value in map

   \Move(0,-25)
   \MarkLoc(tmp)
   \EntryExit(-1,0,1.1,0)
   \Node (tag) (--#2--)
   \ORectNode(ref) (-- --)
   \EntryExit(0,0,0,0)
   \IF \EqInt(#1,1) \THEN		% ordinary ref
      \Move(100,\Val\I)
      \OD(#3, #4, #5, 1)  
   \FI
   \IF \EqInt(#1,2) \THEN		% ref in a set
 	\ODset(#3,#4,#5)
   \FI
   \IF \EqInt(#1,3) \THEN		% key in a map
	    \Move(-160,\Val\I)
	    \Move(0,-10)    
	    \OD(#3, #4, #5,3)  
   \FI
   \IF \EqInt(#1,4) \THEN		% ref to a set
	       \OD(#3,#4,#5,4)
   \FI
   \IF \EqInt(#1,5) \THEN		% ref to a map
	          \ODmap(#3,#4,#5,#6,#7,#8,#1)
   \FI
   \MoveToLoc(tmp)
   \I-100;
   \MarkGLoc(bottom)
\FI
   }
 
%% Draw the object in a rectangular box
%% Show primitive values
%% Show references as arrows to objects
\Define \OD(4)
{  
%% parms:
%%    #1:  class of the object 
%%    #2:  table of primitive fields 
%%    #3:  table of object fields 
%%    #4:  See parm #1 for Obj above

   \SaveAll
%%%%  J is depth of recursion
\J+1;
\IF \GtInt (\J,2) \THEN
      \Move(-400,-200)
      \J=0;
    
\FI
\Move (-20,60)   %%% ???  
\Node(obj)(--\underline{#1}--)
%% topLeft and bottomRight are used to draw the rectangle
\MoveToExit(1.02,0)    \MarkGLoc(maxRight)   %% maxRight stores largest primitive
\MoveToExit(-1.2,1.2)  \MarkLoc(topLeft)
\Move (50,-70)
 
\MinNodeSize(30,15)
\PenSize(0.25pt)	%% Lightweight arrow, might cross over other objects

%% Draw the arrow for a reference
\IF \EqInt (#4,3) \THEN 	%% key in a map
     { \CurvedEdgeAt (ref,-0.8,0, obj,1,0.5)(180,0.6,90, 0.6) }
	\ELSE
	\IF \EqInt (#4,6) \THEN		%% calling OD directly 
		{ } 			%% no arrow
        \ELSE
     	{ \CurvedEdgeAt (ref,0.8,0, obj,-1,0.8)(0,0.6,180, 0.6) }
   	\FI
   \FI

\Move(-45,70)
\PenSize(0.75pt)

{\MarkGLoc (maxL) \Move (0,100) \MarkGLoc(topL) }

\Indirect <#2>(0,9) { \Var }		%% Draw primitive fields
\MarkLoc(bottomTemp)
\MoveToLoc(maxRight) \MarkLoc(maxRightLocal)
\MoveToLoc(bottomTemp)

\Indirect<#3>(0,8) { \Obj   }		%% Draw object fields

%% Set maxRightLocal in case there were no primitive fields
\Move(100,0) \MarkLoc(maxRightTemp)
\MoveToRightmost(maxRightLocal,maxRightTemp)
\MarkLoc(maxRightLocal)
{
%% Find bottom right corner of rectangle
\MoveToLoc(bottom)
\Move(100,0) \MarkLoc(bottomR)
\MoveToLoc(maxRightLocal) \Move(0,-100) \MarkLoc(rightB)
\MoveToLL(bottom,bottomR)(maxRightLocal,rightB)
\Move(10,-20) 
\MarkLoc(bottomRight)
\CSeg \DrawRect(bottomRight,topLeft)
}
   
   \RecallAll
\J-1;
}

%% Draw a ref to a Set
\Define \ObjSet(6)
{	
%% parms:
%%    #1: 0 iff this field has no Object ref
%%    #2: name of the field
%%    #3: class of the field  (for object types)
%%    #4: table of primitive fields in the referred object - wrappers or Strings
%%    #5: table of object fields in the referred object
%%    #6: size of the set 

  % Base case:  #1 = 0

\IF \GtInt (#1, 0) \THEN		% Recursive case
	\MarkLoc(tmp)
   	\Node (tag) (--#2--)
	\Move (50,0)
	\ORectNode(ref) (-- --)
	\Move(50,\Val\I)
	\Move (35,20)
	\ODset(#3, #4, #5,#6)  
	\MoveToLoc(tmp)
	\I-135;
   \FI
}

%% Draw a set with the refs at 'random' places in the box
\Define \ODset(4)    
{  
%% parms:
%%    #1:  class of the object
%%    #2:  table of primitive fields 
%%    #3:  table of object fields 
%%    #4:  size of the set        
   \SaveAll
  
\MinNodeSize (90,110)
\RectNode(obj)(-- #1~~~~~~~~~~~~~~~~~~~~~~~--)
\MinNodeSize (20,10)

{ \PenSize(0.25pt)	%% Lightweight arrow, might cross over other objects
  %% Draw arrow to this Set object
  \CurvedEdgeAt (ref,0.8,0, obj,-1,0.8)(10,0.6,180, 0.6) }
  \PenSize(0.75pt)
   \IntVar \X      
   \IntVar \Y

   \Move(-50,70)
   { \Line(110,0) }
   \Move (15,-15)
   \Node(tag)(--size--)
   \MoveToExit(1,0)
   \Move(25,0)
   \ORectNode(tag)(--#4--)
   \Move(-90,0)
   \MarkLoc(start)

%% Place the fields at seemingly random locations
   \Indirect <#2>(0,9) { 		%% Draw the primitive fields
 	\X+34;
 	\K=\X;
	\K/90;
	\K*90;
 	\X-\K;		%% x + 34 mod 90
	\Y+73;
	\K=\Y;
	\K/110;
	\K*110;
	\Y-\K;		%% y + 73 mod 110
	\Y+20;		%% Allow for size at top
 		\MoveToLoc(start)
		\Move (\Val\X,-\Val\Y)
 		\Var
}

%% Place the fields at seemingly random locations
   \Indirect<#3>(0,9) { 		%% Draw the object fields
 	\X+34;
 	\K=\X;
 	\K/90;
 	\K*90;
 	\X-\K;		%% x + 72 mod 90
	\Y+73;
	\K=\Y;
	\K/110;
	\K*110;
	\Y-\K;		%% y + 73 mod 110
	\Y+20;		%% Allow for size at top
 	\MoveToLoc(start)
	\Move (\Val\X,-\Val\Y)
		 \Obj   
	}
   \RecallAll
}

%% Draw a map object in a rectangle
%% Value refs point to the right
%% Key refs point to the left
%% Note largest value field to draw the rectangle
\Define \ODmap(7)
{  
%% parms:
%%    #1:  class of the object 
%%    #2:  size
%%    #3:  table of primitive keys 
%%    #4:  table of object keys 
%%    #5:  table of primitive values 
%%    #6:  table of object values 
%%    #7:  see parm #1 for Obj above
   \SaveAll
  
\Move(-20,-60) 
\Node(obj)(--\underline{#1}--)
\MoveToExit(1.02,0) 	\MarkLoc (maxRightLocal)
			\MarkLoc (maxRight)
\MoveToExit(-1.2,1.2)\MarkLoc (topLeft)

\MinNodeSize(30,15)
\PenSize(0.25pt)		%% Use Lightweight pen, may cross other objects

\IF \EqInt(#7,3) \THEN
     \CurvedEdgeAt (ref,-0.8,0, obj,1,0.5)(180,0.6,90, 0.6) 
     \ELSE
	\IF \EqInt(#7,6) \THEN		%% calling ODmap directly
	    {  }
	\ELSE
	  \CurvedEdgeAt (ref,0.8,0, obj,-1,0.8)(10,0.6,180, 0.6) 
        \FI
\FI
	
\PenSize(0.75pt)
\Move(20,-35)
   
   %% Show size and headings for keys and values
   { \Node(tag)(--size--) \Move(40,0) \ORectNode(tag)(--#2--) }
   \Move (10,-20)
   { \Node(tag)(--\underline{keys}--) 
     \Move(75,0) \Node(tag)(--\underline {values}--) 
   }
   \Move (-50,0)
   \MarkLoc(keys)

   {\MarkGLoc (maxL) \Move (0,100) \MarkGLoc(topL) }   %% ??

  %% keys 
   \Indirect <#3>(0,9) { \Var  }	%% Draw the primitive keys

   \Indirect<#4>(0,8) { \Obj   }	%% Draw the object keys
   \MarkGLoc(bottom)
\MoveToLoc(keys)
\Move(83, -1)
{
  %% values 
   \Indirect <#5>(0,9) { \Var  }	%% Draw the primitive values
   \MarkLoc(bottomTemp)
   \MoveToLoc(bottomTemp)

   \Indirect<#6>(0,8) { \Obj   }	%% Draw the object values

%% Keys may not be primitive, affects size of box
\Move(100,0)	\MarkLoc(maxRightTemp)
\MoveToRightmost(maxRightLocal,maxRightTemp)
\MarkLoc(maxRightLocal)
}
{
%% Find bottom right corner of rectangle
\MoveToLoc(bottom)
\Move(100,0)	\MarkLoc(bottomR)
\MoveToLoc(maxRightLocal)	\Move(0,-100)	\MarkLoc(rightB)
\MoveToLL(bottom,bottomR)(maxRightLocal,rightB)
\Move(10,-20)
\MarkLoc(bottomRight)
\CSeg \DrawRect(bottomRight,topLeft)
}
   \RecallAll
}

%% Move the pen to the loc which is furthest right
\Define \MoveToRightmost(2)
{
%% #1:  A loc
%% #2:  A loc
\MoveTo(-200,100)	\MarkLoc(upLeft)
\MoveTo(-200,-100)	\MarkLoc(downLeft)
\MoveToLoc(#1)		\Move(100,0)	\MarkLoc(right)
\MoveToLL(upLeft,downLeft)(#1,right)	\MarkLoc(leftA)
\MoveToLoc(#2)		\Move(100,0)	\MarkLoc(right)
\MoveToLL(upLeft,downLeft)(#2,right)	\MarkLoc(leftB)

	\LSeg\Q(leftA,#1)
	\LSeg\R(leftB,#2)
	\MoveToLoc(#1)
	\IF \GtDec (\R,\Q) \THEN
	    \MoveToLoc(#2)
	    \FI
}

\hspace{-4cm}

\Indirect \Table <ObjectPrims>
{  
	1, head, \LLHead  &
    1, tail, \LLTail &
    1, size, \LLSize
}

\Indirect \Table <ObjectDiagram>
 {  0, home,Address,addrAPrims,addrAObjs, , , }
\Scale (1,1)
\OD (\ObjectName, ObjectPrims, ObjectDiagram, 6)

%Arrow from head
%%Draws the arrow from head to first node
\Move(5,43)
\FcNode(Head)
\Move(45,0)
\FcNode(HeadRight)
\Move(0,-80)
\FcNode(HeadDown)
\Move(-100,0)
\FcNode(HeadLeft)
\Move(0,-14)
\FcNode(HeadBottom)
\ArrowHeads(1)
\Edge(HeadLeft,HeadBottom)
\ArrowHeads(0)
\Edge(Head,HeadRight)
\VHEdge(HeadRight,HeadDown)
\Edge(HeadDown,HeadLeft)
\EndDraw

%Node
\begin{tikzpicture}[abstract, rectangle split, rectangle split parts=2, text justified]
\foreach \x in {0}
    \foreach \y in {1,...,\numofnodess} 
       {\node (NODES) [abstract, rectangle split, rectangle split parts=2, text justified]  (\x\y) at (5, 2.9*\y) 
       		{\textbf{Node}
            \nodepart{second}next  \framebox[1.1\width]{arrow}
            \newline prev  \framebox[1.1\width]{arrow}
            \newline value  \framebox[1.1\width]{arrow}
            
            };}
                      
\foreach \x in {0}
    \foreach \y in {1,...,\numofnodess}            
{\node (NODETHING) [abstract, rectangle split, rectangle split parts=2, text justified] (\x\y) at (10, 2.6*\y) 
       		{\textbf{\VTop}
            \nodepart{second}Name  \framebox[1.1\width]{\VName}
            \newline Num  \framebox[1.1\width]{\VNum}
            \newline Address  \framebox[1.1\width]{arrow}};}

\foreach \x in {0}
    \foreach \y in {1,...,\numofnodess} 
{\node (NODETHING2) [abstract, rectangle split, rectangle split parts=2, text justified] (\x\y) at (15, 2.6*\y) 
       		{\textbf{\AName}
            \nodepart{second}Street  \framebox[1.1\width]{\AStreet}
            \newline City  \framebox[1.1\width]{\ACity}
            \newline State  \framebox[1.1\width]{\AState}
            \newline Zip  \framebox[1.1\width]{\AZip}};}           
\end{tikzpicture}

%%Tail
%%Draws the tail node
\hspace{.05cm}
\Draw


\MinNodeSize (20,10)
\ArrowHeads(1)
\ArrowSpec(F,10,4,5)
%% \DecVar \MaxX

%% Display the primitive values in an object
%% Called once for each primitive field
\Define \Var(3)
{	
%% parms:
%%    #1: 0 if this var has no primitive value
%%    #2: name of the var
%%    #3: value in a box
	
   \IF \GtInt (#1,0)  \THEN
	\MarkLoc(left)
 	\Move(0,-27)
	\EntryExit(-1,0,-1,0)
	\Node (tag) (--#2--)
	%% Note the longest field value, to draw the box
	{ 
	\MoveToExit(1,0)
	  \ORectNode (tag) (--#3--)
	\MoveToExit(1,0)  \MarkLoc(right)
	\Move(50,0)  \MarkLoc(rightPlus) 
	\MoveToLL(right,rightPlus)(topL,maxL)
	\MarkGLoc(bottomL)  
	\LSeg\Q(maxL,maxRight)
	\LSeg\R(bottomL,right)
	\IF \GtDec (\R,\Q) \THEN
	    \MoveToLoc(right)
	    \MarkGLoc(maxRight)
	    \FI
	}
	\EntryExit(0,0,0,0)
\MarkGLoc(bottom)
   \FI
}


%% Display the object references in an object
%% Called once for each primitive field
%% Object is a recursive data structure
\Define \Obj(8)
{	
%% parms:
%%    #1: 0 do not display
%%        1 display ordinary ref
%%	  2 display ref in a set
%%        3 display a key in a map
%%        4 display a ref to a set
%%        5 display a ref to a map
%%  	  6 display object only, calling OD directly, no arrow
%%    #2: name of the field
%%    #3: class of the field  (for object types)
%%    #4: table of primitive fields in the referred object
%% 	   or size for a map
%%    #5: table of object fields in the referred object
%%         or table of primitive keys for maps
%%    #6: table of object keys in the referred object
%%    #7: table of primitive values in the referred object
%%    #8: table of object values in the referred object

  % Base case for recursion:  #1 = 0

\IF \GtInt (#1, 0) \THEN		% Recursive case
					% for ordinary field or value in map

   \Move(0,-25)
   \MarkLoc(tmp)
   \EntryExit(-1,0,1.1,0)
   \Node (tag) (--#2--)
   \ORectNode(ref) (-- --)
   \EntryExit(0,0,0,0)
   \IF \EqInt(#1,1) \THEN		% ordinary ref
      \Move(100,\Val\I)
      \OD(#3, #4, #5, 1)  
   \FI
   \IF \EqInt(#1,2) \THEN		% ref in a set
 	\ODset(#3,#4,#5)
   \FI
   \IF \EqInt(#1,3) \THEN		% key in a map
	    \Move(-160,\Val\I)
	    \Move(0,-10)    
	    \OD(#3, #4, #5,3)  
   \FI
   \IF \EqInt(#1,4) \THEN		% ref to a set
	       \OD(#3,#4,#5,4)
   \FI
   \IF \EqInt(#1,5) \THEN		% ref to a map
	          \ODmap(#3,#4,#5,#6,#7,#8,#1)
   \FI
   \MoveToLoc(tmp)
   \I-100;
   \MarkGLoc(bottom)
\FI
   }
 
%% Draw the object in a rectangular box
%% Show primitive values
%% Show references as arrows to objects
\Define \OD(4)
{  
%% parms:
%%    #1:  class of the object 
%%    #2:  table of primitive fields 
%%    #3:  table of object fields 
%%    #4:  See parm #1 for Obj above

   \SaveAll
%%%%  J is depth of recursion
\J+1;
\IF \GtInt (\J,2) \THEN
      \Move(-400,-200)
      \J=0;
    
\FI
\Move (-20,60)   %%% ???  
\Node(obj)(--\underline{#1}--)
%% topLeft and bottomRight are used to draw the rectangle
\MoveToExit(1.02,0)    \MarkGLoc(maxRight)   %% maxRight stores largest primitive
\MoveToExit(-1.2,1.2)  \MarkLoc(topLeft)
\Move (50,-70)
 
\MinNodeSize(30,15)
\PenSize(0.25pt)	%% Lightweight arrow, might cross over other objects

%% Draw the arrow for a reference
\IF \EqInt (#4,3) \THEN 	%% key in a map
     { \CurvedEdgeAt (ref,-0.8,0, obj,1,0.5)(180,0.6,90, 0.6) }
	\ELSE
	\IF \EqInt (#4,6) \THEN		%% calling OD directly 
		{ } 			%% no arrow
        \ELSE
     	{ \CurvedEdgeAt (ref,0.8,0, obj,-1,0.8)(0,0.6,180, 0.6) }
   	\FI
   \FI

\Move(-45,70)
\PenSize(0.75pt)

{\MarkGLoc (maxL) \Move (0,100) \MarkGLoc(topL) }

\Indirect <#2>(0,9) { \Var }		%% Draw primitive fields
\MarkLoc(bottomTemp)
\MoveToLoc(maxRight) \MarkLoc(maxRightLocal)
\MoveToLoc(bottomTemp)

\Indirect<#3>(0,8) { \Obj   }		%% Draw object fields

%% Set maxRightLocal in case there were no primitive fields
\Move(100,0) \MarkLoc(maxRightTemp)
\MoveToRightmost(maxRightLocal,maxRightTemp)
\MarkLoc(maxRightLocal)
{
%% Find bottom right corner of rectangle
\MoveToLoc(bottom)
\Move(100,0) \MarkLoc(bottomR)
\MoveToLoc(maxRightLocal) \Move(0,-100) \MarkLoc(rightB)
\MoveToLL(bottom,bottomR)(maxRightLocal,rightB)
\Move(10,-20) 
\MarkLoc(bottomRight)
\CSeg \DrawRect(bottomRight,topLeft)
}
   
   \RecallAll
\J-1;
}

%% Draw a ref to a Set
\Define \ObjSet(6)
{	
%% parms:
%%    #1: 0 iff this field has no Object ref
%%    #2: name of the field
%%    #3: class of the field  (for object types)
%%    #4: table of primitive fields in the referred object - wrappers or Strings
%%    #5: table of object fields in the referred object
%%    #6: size of the set 

  % Base case:  #1 = 0

\IF \GtInt (#1, 0) \THEN		% Recursive case
	\MarkLoc(tmp)
   	\Node (tag) (--#2--)
	\Move (50,0)
	\ORectNode(ref) (-- --)
	\Move(50,\Val\I)
	\Move (35,20)
	\ODset(#3, #4, #5,#6)  
	\MoveToLoc(tmp)
	\I-135;
   \FI
}

%% Draw a set with the refs at 'random' places in the box
\Define \ODset(4)    
{  
%% parms:
%%    #1:  class of the object
%%    #2:  table of primitive fields 
%%    #3:  table of object fields 
%%    #4:  size of the set        
   \SaveAll
  
\MinNodeSize (90,110)
\RectNode(obj)(-- #1~~~~~~~~~~~~~~~~~~~~~~~--)
\MinNodeSize (20,10)

{ \PenSize(0.25pt)	%% Lightweight arrow, might cross over other objects
  %% Draw arrow to this Set object
  \CurvedEdgeAt (ref,0.8,0, obj,-1,0.8)(10,0.6,180, 0.6) }
  \PenSize(0.75pt)
   \IntVar \X      
   \IntVar \Y

   \Move(-50,70)
   { \Line(110,0) }
   \Move (15,-15)
   \Node(tag)(--size--)
   \MoveToExit(1,0)
   \Move(25,0)
   \ORectNode(tag)(--#4--)
   \Move(-90,0)
   \MarkLoc(start)

%% Place the fields at seemingly random locations
   \Indirect <#2>(0,9) { 		%% Draw the primitive fields
 	\X+34;
 	\K=\X;
	\K/90;
	\K*90;
 	\X-\K;		%% x + 34 mod 90
	\Y+73;
	\K=\Y;
	\K/110;
	\K*110;
	\Y-\K;		%% y + 73 mod 110
	\Y+20;		%% Allow for size at top
 		\MoveToLoc(start)
		\Move (\Val\X,-\Val\Y)
 		\Var
}

%% Place the fields at seemingly random locations
   \Indirect<#3>(0,9) { 		%% Draw the object fields
 	\X+34;
 	\K=\X;
 	\K/90;
 	\K*90;
 	\X-\K;		%% x + 72 mod 90
	\Y+73;
	\K=\Y;
	\K/110;
	\K*110;
	\Y-\K;		%% y + 73 mod 110
	\Y+20;		%% Allow for size at top
 	\MoveToLoc(start)
	\Move (\Val\X,-\Val\Y)
		 \Obj   
	}
   \RecallAll
}

%% Draw a map object in a rectangle
%% Value refs point to the right
%% Key refs point to the left
%% Note largest value field to draw the rectangle
\Define \ODmap(7)
{  
%% parms:
%%    #1:  class of the object 
%%    #2:  size
%%    #3:  table of primitive keys 
%%    #4:  table of object keys 
%%    #5:  table of primitive values 
%%    #6:  table of object values 
%%    #7:  see parm #1 for Obj above
   \SaveAll
  
\Move(-20,-60) 
\Node(obj)(--\underline{#1}--)
\MoveToExit(1.02,0) 	\MarkLoc (maxRightLocal)
			\MarkLoc (maxRight)
\MoveToExit(-1.2,1.2)\MarkLoc (topLeft)

\MinNodeSize(30,15)
\PenSize(0.25pt)		%% Use Lightweight pen, may cross other objects

\IF \EqInt(#7,3) \THEN
     \CurvedEdgeAt (ref,-0.8,0, obj,1,0.5)(180,0.6,90, 0.6) 
     \ELSE
	\IF \EqInt(#7,6) \THEN		%% calling ODmap directly
	    {  }
	\ELSE
	  \CurvedEdgeAt (ref,0.8,0, obj,-1,0.8)(10,0.6,180, 0.6) 
        \FI
\FI
	
\PenSize(0.75pt)
\Move(20,-35)
   
   %% Show size and headings for keys and values
   { \Node(tag)(--size--) \Move(40,0) \ORectNode(tag)(--#2--) }
   \Move (10,-20)
   { \Node(tag)(--\underline{keys}--) 
     \Move(75,0) \Node(tag)(--\underline {values}--) 
   }
   \Move (-50,0)
   \MarkLoc(keys)

   {\MarkGLoc (maxL) \Move (0,100) \MarkGLoc(topL) }   %% ??

  %% keys 
   \Indirect <#3>(0,9) { \Var  }	%% Draw the primitive keys

   \Indirect<#4>(0,8) { \Obj   }	%% Draw the object keys
   \MarkGLoc(bottom)
\MoveToLoc(keys)
\Move(83, -1)
{
  %% values 
   \Indirect <#5>(0,9) { \Var  }	%% Draw the primitive values
   \MarkLoc(bottomTemp)
   \MoveToLoc(bottomTemp)

   \Indirect<#6>(0,8) { \Obj   }	%% Draw the object values

%% Keys may not be primitive, affects size of box
\Move(100,0)	\MarkLoc(maxRightTemp)
\MoveToRightmost(maxRightLocal,maxRightTemp)
\MarkLoc(maxRightLocal)
}
{
%% Find bottom right corner of rectangle
\MoveToLoc(bottom)
\Move(100,0)	\MarkLoc(bottomR)
\MoveToLoc(maxRightLocal)	\Move(0,-100)	\MarkLoc(rightB)
\MoveToLL(bottom,bottomR)(maxRightLocal,rightB)
\Move(10,-20)
\MarkLoc(bottomRight)
\CSeg \DrawRect(bottomRight,topLeft)
}
   \RecallAll
}

%% Move the pen to the loc which is furthest right
\Define \MoveToRightmost(2)
{
%% #1:  A loc
%% #2:  A loc
\MoveTo(-200,100)	\MarkLoc(upLeft)
\MoveTo(-200,-100)	\MarkLoc(downLeft)
\MoveToLoc(#1)		\Move(100,0)	\MarkLoc(right)
\MoveToLL(upLeft,downLeft)(#1,right)	\MarkLoc(leftA)
\MoveToLoc(#2)		\Move(100,0)	\MarkLoc(right)
\MoveToLL(upLeft,downLeft)(#2,right)	\MarkLoc(leftB)

	\LSeg\Q(leftA,#1)
	\LSeg\R(leftB,#2)
	\MoveToLoc(#1)
	\IF \GtDec (\R,\Q) \THEN
	    \MoveToLoc(#2)
	    \FI
}


\Indirect \Table <ObjectPrims>
{  
	1, next, \THead &
    1, prev, \TTail &
    1, value, \TSize
}
\Indirect \Table <ObjectDiagram>
 {  0, home,Address,addrAPrims,addrAObjs, , , }
\Scale (1,1)
\OD (\TName, ObjectPrims, ObjectDiagram, 6)

%Arrow from Tail
%%Draws the arrow from tail node to tail
\Move(48,15)
\FcNode(Tail)
\Move(45,0)
\FcNode(Temp2)
\Move(0,60)
\FcNode(Temp3)
\Move(-100,0)
\FcNode(Temp4)
\Move(0,14)
\FcNode(Temp5)
\ArrowHeads(1)
\Edge(Temp4,Temp5)
\ArrowHeads(0)
\Edge(Tail,Temp2)
\Edge(Temp2,Temp3)
\Edge(Temp3,Temp4)
\EndDraw

\caption {\Caption}

\label {fig:ObjectDiagram}

\end {figure}

%%%%%%%5 input
\newpage

\newcommand{\numofnodesss}{5}


\begin {figure}

%Head
%%Draws the head node
\Draw


\MinNodeSize (20,10)
\ArrowHeads(1)
\ArrowSpec(F,10,4,5)
%% \DecVar \MaxX

%% Display the primitive values in an object
%% Called once for each primitive field
\Define \Var(3)
{	
%% parms:
%%    #1: 0 if this var has no primitive value
%%    #2: name of the var
%%    #3: value in a box
	
   \IF \GtInt (#1,0)  \THEN
	\MarkLoc(left)
 	\Move(0,-27)
	\EntryExit(-1,0,-1,0)
	\Node (tag) (--#2--)
	%% Note the longest field value, to draw the box
	{ 
	\MoveToExit(1,0)
	  \ORectNode (tag) (--#3--)
	\MoveToExit(1,0)  \MarkLoc(right)
	\Move(50,0)  \MarkLoc(rightPlus) 
	\MoveToLL(right,rightPlus)(topL,maxL)
	\MarkGLoc(bottomL)  
	\LSeg\Q(maxL,maxRight)
	\LSeg\R(bottomL,right)
	\IF \GtDec (\R,\Q) \THEN
	    \MoveToLoc(right)
	    \MarkGLoc(maxRight)
	    \FI
	}
	\EntryExit(0,0,0,0)
\MarkGLoc(bottom)
   \FI
}


%% Display the object references in an object
%% Called once for each primitive field
%% Object is a recursive data structure
\Define \Obj(8)
{	
%% parms:
%%    #1: 0 do not display
%%        1 display ordinary ref
%%	  2 display ref in a set
%%        3 display a key in a map
%%        4 display a ref to a set
%%        5 display a ref to a map
%%  	  6 display object only, calling OD directly, no arrow
%%    #2: name of the field
%%    #3: class of the field  (for object types)
%%    #4: table of primitive fields in the referred object
%% 	   or size for a map
%%    #5: table of object fields in the referred object
%%         or table of primitive keys for maps
%%    #6: table of object keys in the referred object
%%    #7: table of primitive values in the referred object
%%    #8: table of object values in the referred object

  % Base case for recursion:  #1 = 0

\IF \GtInt (#1, 0) \THEN		% Recursive case
					% for ordinary field or value in map

   \Move(0,-25)
   \MarkLoc(tmp)
   \EntryExit(-1,0,1.1,0)
   \Node (tag) (--#2--)
   \ORectNode(ref) (-- --)
   \EntryExit(0,0,0,0)
   \IF \EqInt(#1,1) \THEN		% ordinary ref
      \Move(100,\Val\I)
      \OD(#3, #4, #5, 1)  
   \FI
   \IF \EqInt(#1,2) \THEN		% ref in a set
 	\ODset(#3,#4,#5)
   \FI
   \IF \EqInt(#1,3) \THEN		% key in a map
	    \Move(-160,\Val\I)
	    \Move(0,-10)    
	    \OD(#3, #4, #5,3)  
   \FI
   \IF \EqInt(#1,4) \THEN		% ref to a set
	       \OD(#3,#4,#5,4)
   \FI
   \IF \EqInt(#1,5) \THEN		% ref to a map
	          \ODmap(#3,#4,#5,#6,#7,#8,#1)
   \FI
   \MoveToLoc(tmp)
   \I-100;
   \MarkGLoc(bottom)
\FI
   }
 
%% Draw the object in a rectangular box
%% Show primitive values
%% Show references as arrows to objects
\Define \OD(4)
{  
%% parms:
%%    #1:  class of the object 
%%    #2:  table of primitive fields 
%%    #3:  table of object fields 
%%    #4:  See parm #1 for Obj above

   \SaveAll
%%%%  J is depth of recursion
\J+1;
\IF \GtInt (\J,2) \THEN
      \Move(-400,-200)
      \J=0;
    
\FI
\Move (-20,60)   %%% ???  
\Node(obj)(--\underline{#1}--)
%% topLeft and bottomRight are used to draw the rectangle
\MoveToExit(1.02,0)    \MarkGLoc(maxRight)   %% maxRight stores largest primitive
\MoveToExit(-1.2,1.2)  \MarkLoc(topLeft)
\Move (50,-70)
 
\MinNodeSize(30,15)
\PenSize(0.25pt)	%% Lightweight arrow, might cross over other objects

%% Draw the arrow for a reference
\IF \EqInt (#4,3) \THEN 	%% key in a map
     { \CurvedEdgeAt (ref,-0.8,0, obj,1,0.5)(180,0.6,90, 0.6) }
	\ELSE
	\IF \EqInt (#4,6) \THEN		%% calling OD directly 
		{ } 			%% no arrow
        \ELSE
     	{ \CurvedEdgeAt (ref,0.8,0, obj,-1,0.8)(0,0.6,180, 0.6) }
   	\FI
   \FI

\Move(-45,70)
\PenSize(0.75pt)

{\MarkGLoc (maxL) \Move (0,100) \MarkGLoc(topL) }

\Indirect <#2>(0,9) { \Var }		%% Draw primitive fields
\MarkLoc(bottomTemp)
\MoveToLoc(maxRight) \MarkLoc(maxRightLocal)
\MoveToLoc(bottomTemp)

\Indirect<#3>(0,8) { \Obj   }		%% Draw object fields

%% Set maxRightLocal in case there were no primitive fields
\Move(100,0) \MarkLoc(maxRightTemp)
\MoveToRightmost(maxRightLocal,maxRightTemp)
\MarkLoc(maxRightLocal)
{
%% Find bottom right corner of rectangle
\MoveToLoc(bottom)
\Move(100,0) \MarkLoc(bottomR)
\MoveToLoc(maxRightLocal) \Move(0,-100) \MarkLoc(rightB)
\MoveToLL(bottom,bottomR)(maxRightLocal,rightB)
\Move(10,-20) 
\MarkLoc(bottomRight)
\CSeg \DrawRect(bottomRight,topLeft)
}
   
   \RecallAll
\J-1;
}

%% Draw a ref to a Set
\Define \ObjSet(6)
{	
%% parms:
%%    #1: 0 iff this field has no Object ref
%%    #2: name of the field
%%    #3: class of the field  (for object types)
%%    #4: table of primitive fields in the referred object - wrappers or Strings
%%    #5: table of object fields in the referred object
%%    #6: size of the set 

  % Base case:  #1 = 0

\IF \GtInt (#1, 0) \THEN		% Recursive case
	\MarkLoc(tmp)
   	\Node (tag) (--#2--)
	\Move (50,0)
	\ORectNode(ref) (-- --)
	\Move(50,\Val\I)
	\Move (35,20)
	\ODset(#3, #4, #5,#6)  
	\MoveToLoc(tmp)
	\I-135;
   \FI
}

%% Draw a set with the refs at 'random' places in the box
\Define \ODset(4)    
{  
%% parms:
%%    #1:  class of the object
%%    #2:  table of primitive fields 
%%    #3:  table of object fields 
%%    #4:  size of the set        
   \SaveAll
  
\MinNodeSize (90,110)
\RectNode(obj)(-- #1~~~~~~~~~~~~~~~~~~~~~~~--)
\MinNodeSize (20,10)

{ \PenSize(0.25pt)	%% Lightweight arrow, might cross over other objects
  %% Draw arrow to this Set object
  \CurvedEdgeAt (ref,0.8,0, obj,-1,0.8)(10,0.6,180, 0.6) }
  \PenSize(0.75pt)
   \IntVar \X      
   \IntVar \Y

   \Move(-50,70)
   { \Line(110,0) }
   \Move (15,-15)
   \Node(tag)(--size--)
   \MoveToExit(1,0)
   \Move(25,0)
   \ORectNode(tag)(--#4--)
   \Move(-90,0)
   \MarkLoc(start)

%% Place the fields at seemingly random locations
   \Indirect <#2>(0,9) { 		%% Draw the primitive fields
 	\X+34;
 	\K=\X;
	\K/90;
	\K*90;
 	\X-\K;		%% x + 34 mod 90
	\Y+73;
	\K=\Y;
	\K/110;
	\K*110;
	\Y-\K;		%% y + 73 mod 110
	\Y+20;		%% Allow for size at top
 		\MoveToLoc(start)
		\Move (\Val\X,-\Val\Y)
 		\Var
}

%% Place the fields at seemingly random locations
   \Indirect<#3>(0,9) { 		%% Draw the object fields
 	\X+34;
 	\K=\X;
 	\K/90;
 	\K*90;
 	\X-\K;		%% x + 72 mod 90
	\Y+73;
	\K=\Y;
	\K/110;
	\K*110;
	\Y-\K;		%% y + 73 mod 110
	\Y+20;		%% Allow for size at top
 	\MoveToLoc(start)
	\Move (\Val\X,-\Val\Y)
		 \Obj   
	}
   \RecallAll
}

%% Draw a map object in a rectangle
%% Value refs point to the right
%% Key refs point to the left
%% Note largest value field to draw the rectangle
\Define \ODmap(7)
{  
%% parms:
%%    #1:  class of the object 
%%    #2:  size
%%    #3:  table of primitive keys 
%%    #4:  table of object keys 
%%    #5:  table of primitive values 
%%    #6:  table of object values 
%%    #7:  see parm #1 for Obj above
   \SaveAll
  
\Move(-20,-60) 
\Node(obj)(--\underline{#1}--)
\MoveToExit(1.02,0) 	\MarkLoc (maxRightLocal)
			\MarkLoc (maxRight)
\MoveToExit(-1.2,1.2)\MarkLoc (topLeft)

\MinNodeSize(30,15)
\PenSize(0.25pt)		%% Use Lightweight pen, may cross other objects

\IF \EqInt(#7,3) \THEN
     \CurvedEdgeAt (ref,-0.8,0, obj,1,0.5)(180,0.6,90, 0.6) 
     \ELSE
	\IF \EqInt(#7,6) \THEN		%% calling ODmap directly
	    {  }
	\ELSE
	  \CurvedEdgeAt (ref,0.8,0, obj,-1,0.8)(10,0.6,180, 0.6) 
        \FI
\FI
	
\PenSize(0.75pt)
\Move(20,-35)
   
   %% Show size and headings for keys and values
   { \Node(tag)(--size--) \Move(40,0) \ORectNode(tag)(--#2--) }
   \Move (10,-20)
   { \Node(tag)(--\underline{keys}--) 
     \Move(75,0) \Node(tag)(--\underline {values}--) 
   }
   \Move (-50,0)
   \MarkLoc(keys)

   {\MarkGLoc (maxL) \Move (0,100) \MarkGLoc(topL) }   %% ??

  %% keys 
   \Indirect <#3>(0,9) { \Var  }	%% Draw the primitive keys

   \Indirect<#4>(0,8) { \Obj   }	%% Draw the object keys
   \MarkGLoc(bottom)
\MoveToLoc(keys)
\Move(83, -1)
{
  %% values 
   \Indirect <#5>(0,9) { \Var  }	%% Draw the primitive values
   \MarkLoc(bottomTemp)
   \MoveToLoc(bottomTemp)

   \Indirect<#6>(0,8) { \Obj   }	%% Draw the object values

%% Keys may not be primitive, affects size of box
\Move(100,0)	\MarkLoc(maxRightTemp)
\MoveToRightmost(maxRightLocal,maxRightTemp)
\MarkLoc(maxRightLocal)
}
{
%% Find bottom right corner of rectangle
\MoveToLoc(bottom)
\Move(100,0)	\MarkLoc(bottomR)
\MoveToLoc(maxRightLocal)	\Move(0,-100)	\MarkLoc(rightB)
\MoveToLL(bottom,bottomR)(maxRightLocal,rightB)
\Move(10,-20)
\MarkLoc(bottomRight)
\CSeg \DrawRect(bottomRight,topLeft)
}
   \RecallAll
}

%% Move the pen to the loc which is furthest right
\Define \MoveToRightmost(2)
{
%% #1:  A loc
%% #2:  A loc
\MoveTo(-200,100)	\MarkLoc(upLeft)
\MoveTo(-200,-100)	\MarkLoc(downLeft)
\MoveToLoc(#1)		\Move(100,0)	\MarkLoc(right)
\MoveToLL(upLeft,downLeft)(#1,right)	\MarkLoc(leftA)
\MoveToLoc(#2)		\Move(100,0)	\MarkLoc(right)
\MoveToLL(upLeft,downLeft)(#2,right)	\MarkLoc(leftB)

	\LSeg\Q(leftA,#1)
	\LSeg\R(leftB,#2)
	\MoveToLoc(#1)
	\IF \GtDec (\R,\Q) \THEN
	    \MoveToLoc(#2)
	    \FI
}

\hspace{-4cm}

\Indirect \Table <ObjectPrims>
{  
	1, head, \LLHead  &
    1, tail, \LLTail &
    1, size, \LLSize
}

\Indirect \Table <ObjectDiagram>
 {  0, home,Address,addrAPrims,addrAObjs, , , }
\Scale (1,1)
\OD (\ObjectName, ObjectPrims, ObjectDiagram, 6)

%Arrow from head
%%Draws the arrow from head to first node
\Move(5,43)
\FcNode(Head)
\Move(45,0)
\FcNode(HeadRight)
\Move(0,-80)
\FcNode(HeadDown)
\Move(-100,0)
\FcNode(HeadLeft)
\Move(0,-14)
\FcNode(HeadBottom)
\ArrowHeads(1)
\Edge(HeadLeft,HeadBottom)
\ArrowHeads(0)
\Edge(Head,HeadRight)
\VHEdge(HeadRight,HeadDown)
\Edge(HeadDown,HeadLeft)
\EndDraw

%Node
\begin{tikzpicture}[abstract, rectangle split, rectangle split parts=2, text justified]
\foreach \x in {0}
    \foreach \y in {1,...,\numofnodesss} 
       {\node (NODES) [abstract, rectangle split, rectangle split parts=2, text justified]  (\x\y) at (5, 2.9*\y) 
       		{\textbf{Node}
            \nodepart{second}next  \framebox[1.1\width]{arrow}
            \newline prev  \framebox[1.1\width]{arrow}
            \newline value  \framebox[1.1\width]{arrow}
            
            };}
                      
\foreach \x in {0}
    \foreach \y in {1,...,\numofnodesss}            
{\node (NODETHING) [abstract, rectangle split, rectangle split parts=2, text justified] (\x\y) at (10, 2.6*\y) 
       		{\textbf{\VTop}
            \nodepart{second}Name  \framebox[1.1\width]{\VName}
            \newline Num  \framebox[1.1\width]{\VNum}
            \newline Address  \framebox[1.1\width]{arrow}};}

\foreach \x in {0}
    \foreach \y in {1,...,\numofnodesss} 
{\node (NODETHING2) [abstract, rectangle split, rectangle split parts=2, text justified] (\x\y) at (15, 2.6*\y) 
       		{\textbf{\AName}
            \nodepart{second}Street  \framebox[1.1\width]{\AStreet}
            \newline City  \framebox[1.1\width]{\ACity}
            \newline State  \framebox[1.1\width]{\AState}
            \newline Zip  \framebox[1.1\width]{\AZip}};}           
\end{tikzpicture}

%%Tail
%%Draws the tail node
\hspace{.05cm}
\Draw


\MinNodeSize (20,10)
\ArrowHeads(1)
\ArrowSpec(F,10,4,5)
%% \DecVar \MaxX

%% Display the primitive values in an object
%% Called once for each primitive field
\Define \Var(3)
{	
%% parms:
%%    #1: 0 if this var has no primitive value
%%    #2: name of the var
%%    #3: value in a box
	
   \IF \GtInt (#1,0)  \THEN
	\MarkLoc(left)
 	\Move(0,-27)
	\EntryExit(-1,0,-1,0)
	\Node (tag) (--#2--)
	%% Note the longest field value, to draw the box
	{ 
	\MoveToExit(1,0)
	  \ORectNode (tag) (--#3--)
	\MoveToExit(1,0)  \MarkLoc(right)
	\Move(50,0)  \MarkLoc(rightPlus) 
	\MoveToLL(right,rightPlus)(topL,maxL)
	\MarkGLoc(bottomL)  
	\LSeg\Q(maxL,maxRight)
	\LSeg\R(bottomL,right)
	\IF \GtDec (\R,\Q) \THEN
	    \MoveToLoc(right)
	    \MarkGLoc(maxRight)
	    \FI
	}
	\EntryExit(0,0,0,0)
\MarkGLoc(bottom)
   \FI
}


%% Display the object references in an object
%% Called once for each primitive field
%% Object is a recursive data structure
\Define \Obj(8)
{	
%% parms:
%%    #1: 0 do not display
%%        1 display ordinary ref
%%	  2 display ref in a set
%%        3 display a key in a map
%%        4 display a ref to a set
%%        5 display a ref to a map
%%  	  6 display object only, calling OD directly, no arrow
%%    #2: name of the field
%%    #3: class of the field  (for object types)
%%    #4: table of primitive fields in the referred object
%% 	   or size for a map
%%    #5: table of object fields in the referred object
%%         or table of primitive keys for maps
%%    #6: table of object keys in the referred object
%%    #7: table of primitive values in the referred object
%%    #8: table of object values in the referred object

  % Base case for recursion:  #1 = 0

\IF \GtInt (#1, 0) \THEN		% Recursive case
					% for ordinary field or value in map

   \Move(0,-25)
   \MarkLoc(tmp)
   \EntryExit(-1,0,1.1,0)
   \Node (tag) (--#2--)
   \ORectNode(ref) (-- --)
   \EntryExit(0,0,0,0)
   \IF \EqInt(#1,1) \THEN		% ordinary ref
      \Move(100,\Val\I)
      \OD(#3, #4, #5, 1)  
   \FI
   \IF \EqInt(#1,2) \THEN		% ref in a set
 	\ODset(#3,#4,#5)
   \FI
   \IF \EqInt(#1,3) \THEN		% key in a map
	    \Move(-160,\Val\I)
	    \Move(0,-10)    
	    \OD(#3, #4, #5,3)  
   \FI
   \IF \EqInt(#1,4) \THEN		% ref to a set
	       \OD(#3,#4,#5,4)
   \FI
   \IF \EqInt(#1,5) \THEN		% ref to a map
	          \ODmap(#3,#4,#5,#6,#7,#8,#1)
   \FI
   \MoveToLoc(tmp)
   \I-100;
   \MarkGLoc(bottom)
\FI
   }
 
%% Draw the object in a rectangular box
%% Show primitive values
%% Show references as arrows to objects
\Define \OD(4)
{  
%% parms:
%%    #1:  class of the object 
%%    #2:  table of primitive fields 
%%    #3:  table of object fields 
%%    #4:  See parm #1 for Obj above

   \SaveAll
%%%%  J is depth of recursion
\J+1;
\IF \GtInt (\J,2) \THEN
      \Move(-400,-200)
      \J=0;
    
\FI
\Move (-20,60)   %%% ???  
\Node(obj)(--\underline{#1}--)
%% topLeft and bottomRight are used to draw the rectangle
\MoveToExit(1.02,0)    \MarkGLoc(maxRight)   %% maxRight stores largest primitive
\MoveToExit(-1.2,1.2)  \MarkLoc(topLeft)
\Move (50,-70)
 
\MinNodeSize(30,15)
\PenSize(0.25pt)	%% Lightweight arrow, might cross over other objects

%% Draw the arrow for a reference
\IF \EqInt (#4,3) \THEN 	%% key in a map
     { \CurvedEdgeAt (ref,-0.8,0, obj,1,0.5)(180,0.6,90, 0.6) }
	\ELSE
	\IF \EqInt (#4,6) \THEN		%% calling OD directly 
		{ } 			%% no arrow
        \ELSE
     	{ \CurvedEdgeAt (ref,0.8,0, obj,-1,0.8)(0,0.6,180, 0.6) }
   	\FI
   \FI

\Move(-45,70)
\PenSize(0.75pt)

{\MarkGLoc (maxL) \Move (0,100) \MarkGLoc(topL) }

\Indirect <#2>(0,9) { \Var }		%% Draw primitive fields
\MarkLoc(bottomTemp)
\MoveToLoc(maxRight) \MarkLoc(maxRightLocal)
\MoveToLoc(bottomTemp)

\Indirect<#3>(0,8) { \Obj   }		%% Draw object fields

%% Set maxRightLocal in case there were no primitive fields
\Move(100,0) \MarkLoc(maxRightTemp)
\MoveToRightmost(maxRightLocal,maxRightTemp)
\MarkLoc(maxRightLocal)
{
%% Find bottom right corner of rectangle
\MoveToLoc(bottom)
\Move(100,0) \MarkLoc(bottomR)
\MoveToLoc(maxRightLocal) \Move(0,-100) \MarkLoc(rightB)
\MoveToLL(bottom,bottomR)(maxRightLocal,rightB)
\Move(10,-20) 
\MarkLoc(bottomRight)
\CSeg \DrawRect(bottomRight,topLeft)
}
   
   \RecallAll
\J-1;
}

%% Draw a ref to a Set
\Define \ObjSet(6)
{	
%% parms:
%%    #1: 0 iff this field has no Object ref
%%    #2: name of the field
%%    #3: class of the field  (for object types)
%%    #4: table of primitive fields in the referred object - wrappers or Strings
%%    #5: table of object fields in the referred object
%%    #6: size of the set 

  % Base case:  #1 = 0

\IF \GtInt (#1, 0) \THEN		% Recursive case
	\MarkLoc(tmp)
   	\Node (tag) (--#2--)
	\Move (50,0)
	\ORectNode(ref) (-- --)
	\Move(50,\Val\I)
	\Move (35,20)
	\ODset(#3, #4, #5,#6)  
	\MoveToLoc(tmp)
	\I-135;
   \FI
}

%% Draw a set with the refs at 'random' places in the box
\Define \ODset(4)    
{  
%% parms:
%%    #1:  class of the object
%%    #2:  table of primitive fields 
%%    #3:  table of object fields 
%%    #4:  size of the set        
   \SaveAll
  
\MinNodeSize (90,110)
\RectNode(obj)(-- #1~~~~~~~~~~~~~~~~~~~~~~~--)
\MinNodeSize (20,10)

{ \PenSize(0.25pt)	%% Lightweight arrow, might cross over other objects
  %% Draw arrow to this Set object
  \CurvedEdgeAt (ref,0.8,0, obj,-1,0.8)(10,0.6,180, 0.6) }
  \PenSize(0.75pt)
   \IntVar \X      
   \IntVar \Y

   \Move(-50,70)
   { \Line(110,0) }
   \Move (15,-15)
   \Node(tag)(--size--)
   \MoveToExit(1,0)
   \Move(25,0)
   \ORectNode(tag)(--#4--)
   \Move(-90,0)
   \MarkLoc(start)

%% Place the fields at seemingly random locations
   \Indirect <#2>(0,9) { 		%% Draw the primitive fields
 	\X+34;
 	\K=\X;
	\K/90;
	\K*90;
 	\X-\K;		%% x + 34 mod 90
	\Y+73;
	\K=\Y;
	\K/110;
	\K*110;
	\Y-\K;		%% y + 73 mod 110
	\Y+20;		%% Allow for size at top
 		\MoveToLoc(start)
		\Move (\Val\X,-\Val\Y)
 		\Var
}

%% Place the fields at seemingly random locations
   \Indirect<#3>(0,9) { 		%% Draw the object fields
 	\X+34;
 	\K=\X;
 	\K/90;
 	\K*90;
 	\X-\K;		%% x + 72 mod 90
	\Y+73;
	\K=\Y;
	\K/110;
	\K*110;
	\Y-\K;		%% y + 73 mod 110
	\Y+20;		%% Allow for size at top
 	\MoveToLoc(start)
	\Move (\Val\X,-\Val\Y)
		 \Obj   
	}
   \RecallAll
}

%% Draw a map object in a rectangle
%% Value refs point to the right
%% Key refs point to the left
%% Note largest value field to draw the rectangle
\Define \ODmap(7)
{  
%% parms:
%%    #1:  class of the object 
%%    #2:  size
%%    #3:  table of primitive keys 
%%    #4:  table of object keys 
%%    #5:  table of primitive values 
%%    #6:  table of object values 
%%    #7:  see parm #1 for Obj above
   \SaveAll
  
\Move(-20,-60) 
\Node(obj)(--\underline{#1}--)
\MoveToExit(1.02,0) 	\MarkLoc (maxRightLocal)
			\MarkLoc (maxRight)
\MoveToExit(-1.2,1.2)\MarkLoc (topLeft)

\MinNodeSize(30,15)
\PenSize(0.25pt)		%% Use Lightweight pen, may cross other objects

\IF \EqInt(#7,3) \THEN
     \CurvedEdgeAt (ref,-0.8,0, obj,1,0.5)(180,0.6,90, 0.6) 
     \ELSE
	\IF \EqInt(#7,6) \THEN		%% calling ODmap directly
	    {  }
	\ELSE
	  \CurvedEdgeAt (ref,0.8,0, obj,-1,0.8)(10,0.6,180, 0.6) 
        \FI
\FI
	
\PenSize(0.75pt)
\Move(20,-35)
   
   %% Show size and headings for keys and values
   { \Node(tag)(--size--) \Move(40,0) \ORectNode(tag)(--#2--) }
   \Move (10,-20)
   { \Node(tag)(--\underline{keys}--) 
     \Move(75,0) \Node(tag)(--\underline {values}--) 
   }
   \Move (-50,0)
   \MarkLoc(keys)

   {\MarkGLoc (maxL) \Move (0,100) \MarkGLoc(topL) }   %% ??

  %% keys 
   \Indirect <#3>(0,9) { \Var  }	%% Draw the primitive keys

   \Indirect<#4>(0,8) { \Obj   }	%% Draw the object keys
   \MarkGLoc(bottom)
\MoveToLoc(keys)
\Move(83, -1)
{
  %% values 
   \Indirect <#5>(0,9) { \Var  }	%% Draw the primitive values
   \MarkLoc(bottomTemp)
   \MoveToLoc(bottomTemp)

   \Indirect<#6>(0,8) { \Obj   }	%% Draw the object values

%% Keys may not be primitive, affects size of box
\Move(100,0)	\MarkLoc(maxRightTemp)
\MoveToRightmost(maxRightLocal,maxRightTemp)
\MarkLoc(maxRightLocal)
}
{
%% Find bottom right corner of rectangle
\MoveToLoc(bottom)
\Move(100,0)	\MarkLoc(bottomR)
\MoveToLoc(maxRightLocal)	\Move(0,-100)	\MarkLoc(rightB)
\MoveToLL(bottom,bottomR)(maxRightLocal,rightB)
\Move(10,-20)
\MarkLoc(bottomRight)
\CSeg \DrawRect(bottomRight,topLeft)
}
   \RecallAll
}

%% Move the pen to the loc which is furthest right
\Define \MoveToRightmost(2)
{
%% #1:  A loc
%% #2:  A loc
\MoveTo(-200,100)	\MarkLoc(upLeft)
\MoveTo(-200,-100)	\MarkLoc(downLeft)
\MoveToLoc(#1)		\Move(100,0)	\MarkLoc(right)
\MoveToLL(upLeft,downLeft)(#1,right)	\MarkLoc(leftA)
\MoveToLoc(#2)		\Move(100,0)	\MarkLoc(right)
\MoveToLL(upLeft,downLeft)(#2,right)	\MarkLoc(leftB)

	\LSeg\Q(leftA,#1)
	\LSeg\R(leftB,#2)
	\MoveToLoc(#1)
	\IF \GtDec (\R,\Q) \THEN
	    \MoveToLoc(#2)
	    \FI
}


\Indirect \Table <ObjectPrims>
{  
	1, next, \THead &
    1, prev, \TTail &
    1, value, \TSize
}
\Indirect \Table <ObjectDiagram>
 {  0, home,Address,addrAPrims,addrAObjs, , , }
\Scale (1,1)
\OD (\TName, ObjectPrims, ObjectDiagram, 6)

%Arrow from Tail
%%Draws the arrow from tail node to tail
\Move(48,15)
\FcNode(Tail)
\Move(45,0)
\FcNode(Temp2)
\Move(0,60)
\FcNode(Temp3)
\Move(-100,0)
\FcNode(Temp4)
\Move(0,14)
\FcNode(Temp5)
\ArrowHeads(1)
\Edge(Temp4,Temp5)
\ArrowHeads(0)
\Edge(Tail,Temp2)
\Edge(Temp2,Temp3)
\Edge(Temp3,Temp4)
\EndDraw

\caption {\Caption}

\label {fig:ObjectDiagram}

\end {figure}


\end{document}